%!TEX root = ../main.tex
\chapter{Podsumowanie}
\label{ch:podsumowanie}

\section{Realizacja celów pracy}
\label{sec:realizacja-celow}

Głównym celem pracy było zaprojektowanie, implementacja oraz analiza bezpieczeństwa zdecentralizowanej infrastruktury klucza publicznego opartej na technologii blockchain, stanowiącej alternatywę dla scentralizowanego modelu PKI X.509 w zakresie Web PKI. Cel ten zrealizowano w pełnym zakresie: powstał działający prototyp - rejestr wiązań domena--klucz utrzymywany przez sieć walidatorów z konsensusem BFT (rozdziały~\ref{ch:projekt}--\ref{ch:implementacja}) - który poddano systematycznej analizie bezpieczeństwa (rozdział~\ref{ch:analiza-bezpieczenstwa}) oraz ocenie ilościowej (rozdział~\ref{ch:eksperymenty}). Poniżej odniesiono uzyskane wyniki do tezy głównej oraz czterech tez pomocniczych.

\paragraph{Teza główna.} Teza głosząca, że DPKI oparta na blockchain z konsensusem BFT może stanowić funkcjonalnie kompletną alternatywę dla klasycznego modelu PKI X.509 w zakresie Web PKI, eliminując pojedynczy punkt zaufania w postaci centrów certyfikacji, została potwierdzona. Zaimplementowany system realizuje pełny cykl życia wiązania domena--klucz - rejestrację, rotację, unieważnienie oraz weryfikację z dowodem (sekcja~\ref{sec:protokoly}) - a więc jest funkcjonalnie kompletny w rozważanym zakresie. Pojedynczy punkt zaufania został wyeliminowany: jak wykazano w sekcji~\ref{sec:porownanie-bezpieczenstwa}, kompromitacja mniejszości walidatorów nie narusza integralności rejestru, a osiągnięcie progu dwóch trzecich prowadzi co najwyżej do wstrzymania systemu, nie zaś do wystawienia fałszywego wiązania - w przeciwieństwie do modelu klasycznego, w którym kompromitacja dowolnego z setek urzędów wystarcza do wystawienia fałszywego certyfikatu dla dowolnej domeny.

\paragraph{Teza pomocnicza pierwsza.} Potwierdzono, że zastąpienie hierarchii CA rozproszonym rejestrem eliminuje klasę ataków polegających na kompromitacji pojedynczego urzędu, przy czym część zagrożeń pozostaje aktualna. Analiza bezpieczeństwa (rozdział~\ref{ch:analiza-bezpieczenstwa}) wykazała odporność na nieautoryzowane wystawienie i modyfikację wiązań, lecz zarazem uczciwie wskazała zagrożenia rezydualne (sekcja~\ref{sec:ograniczenia}): kompromitację DNS w oknie rejestracji oraz utratę klucza prywatnego właściciela. Zgodnie z przewidywaniem tezy, zagrożenia te dotyczą bezpieczeństwa klucza właściciela i integralności procesu weryfikacji, a więc wymagają niezależnego adresowania, którego sam rejestr nie zapewnia.

\paragraph{Teza pomocnicza druga.} Potwierdzono, że wprowadzenie warstwy konsensusu i sieci P2P rodzi nowe wektory ataku, których realizacja jest jednak istotnie trudniejsza niż kompromitacja pojedynczego CA. Przeanalizowane ataki na konsensus i warstwę sieciową (sekcje~\ref{subsec:kompromitacja-walidatora}--\ref{subsec:sybil-eclipse}) wymagają bądź kontroli ponad jednej trzeciej mocy głosów walidatorów, bądź aktywnej manipulacji ruchem sieciowym - w obu przypadkach progu nieporównanie wyższego niż kompromitacja jednego z setek urzędów w modelu klasycznym. Empirycznie potwierdzono przy tym, że awaria mniejszości walidatorów nie narusza bezpieczeństwa systemu, a jedynie - po przekroczeniu progu kworum - jego dostępność.

\paragraph{Teza pomocnicza trzecia.} Potwierdzono, że zapis stanu w drzewie Merkle umożliwia protokół wiązania pozbawiony zaufanego pośrednika, z gwarancjami silniejszymi niż CRL/OCSP przy niższym koszcie weryfikującego. Weryfikacja istnienia, ważności i unieważnienia odbywa się przez sprawdzenie dowodu względem skrótu stanu, bez odpytywania scentralizowanej infrastruktury statusu (sekcja~\ref{sec:api}). Pomiary (sekcja~\ref{subsec:dowody}) wykazały, że dowód ma rozmiar poniżej 700~bajtów i weryfikuje się w czasie około 28~mikrosekund - a więc taniej niż łańcuch certyfikatów X.509 wraz z odpowiedzią OCSP - a jego rozmiar rośnie jedynie logarytmicznie z liczbą domen. Co więcej, system dostarcza weryfikowalnego dowodu nieistnienia wiązania, niedostępnego w modelu klasycznym.

\paragraph{Teza pomocnicza czwarta.} Potwierdzono, że wykorzystanie gotowej warstwy konsensusu CometBFT oraz interfejsu ABCI pozwala zaimplementować funkcjonalną DPKI wyłącznie w warstwie aplikacyjnej. Cała logika rejestru zamknęła się w aplikacji ABCI (sekcja~\ref{sec:aplikacja-abci}), bez modyfikowania protokołu konsensusu, sieci P2P czy mechanizmu replikacji bloków - co potwierdza praktyczną wykonalność budowy DPKI z gotowych komponentów i stanowi zarazem odpowiedź na pytanie o praktyczną realizowalność takiego systemu.

\section{Wnioski}
\label{sec:wnioski}

Z przeprowadzonej pracy płynie wniosek nadrzędny: zdecentralizowana infrastruktura klucza publicznego oparta na konsensusie BFT jest praktycznie wykonalna i adresuje zasadniczą, strukturalną słabość klasycznej PKI - istnienie pojedynczego punktu zaufania o nieograniczonym zakresie. Zamiast bezwarunkowego zaufania do każdego z setek urzędów certyfikacji, system opiera zaufanie na założeniu o uczciwości ponad dwóch trzecich znanego zbioru walidatorów, czyniąc kompromitację mniejszości nieszkodliwą, a kompromitację większości - prowadzącą do wstrzymania, nie zaś do fałszerstwa.

Praca pozwala też sformułować wniosek o naturze kosztu tej zmiany. Bezpieczeństwo i transparentność uzyskuje się tu przez redundancję i jawność, a nie przez zaufanie do pojedynczego podmiotu - a ceną jest replikacja stanu na wszystkich walidatorach (stały narzut pamięciowy na domenę, sekcja~\ref{subsec:wzrost-stanu}) oraz ograniczenie skali zbioru walidatorów wynikające ze złożoności komunikacji konsensusu (sekcja~\ref{subsec:skalowalnosc}). Jest to kompromis odmienny niż w modelu klasycznym, lecz - jak wykazały pomiary - mieszczący się w granicach praktycznej użyteczności dla rozważanego zastosowania.

Na poziomie szczegółowym praca uwidoczniła kilka rozstrzygnięć projektowych o znaczeniu wykraczającym poza sam prototyp. Najważniejszym z nich jest rola granicy determinizmu wyznaczonej przez interfejs ABCI: konieczność utrzymania w pełni deterministycznego wykonania w ramach konsensusu wymusiła umieszczenie niedeterministycznej weryfikacji DNS jako bramki przedkonsensusowej (sekcja~\ref{sec:wlasnosc-domeny}), co okazało się rozwiązaniem godzącym wymóg weryfikacji własności domeny z wymogiem powtarzalności stanu, a zarazem umożliwiającym odtwarzanie historii przy odzyskiwaniu po awarii. Drugim jest trwałość unieważnienia, która - choć pozornie jest jedynie szczegółem implementacyjnym - okazała się kluczowa dla odporności na atak ,,od-unieważnienia'' w scenariuszu kompromitacji klucza (sekcja~\ref{subsec:ataki-kryptograficzne}).

Należy przy tym zachować właściwą miarę w ocenie wyników. Prototyp nie jest gotowym do wdrożenia produktem, a przeprowadzone testy bezpieczeństwa dowodzą odporności na konkretne, rozważone ataki, nie zaś bezpieczeństwa systemu w sensie absolutnym (sekcja~\ref{sec:model-zagrozen}). Podobnie pomiary wydajnościowe, uzyskane w środowisku jednomaszynowym, dają wartości orientacyjne; wiarygodne pozostają natomiast zaobserwowane relacje asymptotyczne, niezależne od środowiska.

\section{Kierunki dalszych prac}
\label{sec:dalsze-prace}

Praca otwiera kilka naturalnych kierunków kontynuacji, sygnalizowanych w poprzednich rozdziałach. Wynikają one bądź z ograniczeń prototypu, bądź z możliwości jego rozszerzenia o funkcje istotne dla wdrożenia produkcyjnego.

\begin{itemize}
    \item \textbf{Trwały magazyn stanu} - zastąpienie magazynu w pamięci operacyjnej wbudowaną bazą klucz--wartość, za istniejącą już abstrakcją interfejsu magazynu (sekcja~\ref{subsec:zarzadzanie-stanem}). Wyeliminowałoby to konieczność odtwarzania całej historii po restarcie węzła, której koszt rośnie liniowo z długością łańcucha.
    \item \textbf{Kompresja ścieżek drzewa} - zastosowanie węzłów rozszerzających, na wzór drzew stosowanych w systemach produkcyjnych, w celu zmniejszenia stałego narzutu pamięciowego na domenę o jeden do dwóch rzędów wielkości (sekcja~\ref{subsec:wzrost-stanu}), bez zmiany liniowej złożoności asymptotycznej.
    \item \textbf{Pełny lekki klient} - rozszerzenie klienta o weryfikację podpisów walidatorów nad nagłówkami bloków, co domknęłoby lukę w scenariuszu eclipse klienta (sekcja~\ref{subsec:sybil-eclipse}) i zapewniło pełną minimalizację zaufania także w warunkach izolacji sieciowej.
    \item \textbf{Skalowanie poziome} - podział rejestru (np.\ według domen najwyższego poziomu) w celu obejścia ograniczenia skali wynikającego z kwadratowej złożoności komunikacji pojedynczego zbioru walidatorów (sekcja~\ref{subsec:skalowalnosc}).
    \item \textbf{Wieloalgorytmiczność} - rozszerzenie obsługi o dodatkowe algorytmy podpisu poza Ed25519, na które schemat danych już pozostawia miejsce (sekcja~\ref{subsec:reprezentacja-cert}), oraz rozważenie algorytmów odpornych na komputery kwantowe.
    \item \textbf{Integracja z istniejącym ekosystemem} - zbadanie ścieżek współdziałania z obecną Web PKI, w szczególności sposobów wykorzystania zweryfikowanego wiązania przez klientów TLS bez konieczności zmiany całego stosu zaufania naraz.
\end{itemize}

Realizacja powyższych kierunków przybliżyłaby prototyp do rozwiązania zdatnego do wdrożenia, nie zmieniając jednak zasadniczego wyniku niniejszej pracy: wykazania, że zdecentralizowana infrastruktura klucza publicznego oparta na konsensusie BFT jest praktycznie wykonalna i stanowi wiarygodną odpowiedź na strukturalne słabości scentralizowanego modelu zaufania.
